% Part: computability
% Chapter: recursive-functions
% Section: general-recursive-functions

\documentclass[../../../include/open-logic-section]{subfiles}

\begin{document}

\olfileid{cmp}{rec}{gen}
\olsection{توابع بازگشتی عمومی}

راه دیگری نیز برای به‌دست‌آوردن مجموعه‌ای از توابع تام وجود دارد. تابع
تام~$f(x,\vec z)$ را 
\emph{منظم} می‌نامیم اگر برای هر دنبالهٔ~$\vec z$ از اعداد طبیعی،
$x$ای وجود داشته باشد چنان‌که $f(x,\vec z)=0$. به بیان دیگر، توابع
منظم دقیقاً توابعی‌اند که می‌توان جست‌وجوی نامحدود را بر آن‌ها اعمال
کرد و سرانجام تابعی تام به دست آورد. می‌توان، بدون ازدست‌دادن هیچ
تابعی، جست‌وجوی نامحدود را به توابع منظم محدود کرد:

\begin{defn}
\ollabel{defn:general-recursive}
مجموعهٔ 
\emph{توابع بازگشتی عمومی} کوچک‌ترین مجموعه از توابع با موضعیت‌های
گوناگون از اعداد طبیعی به اعداد طبیعی است که صفر، جانشین و افکنش‌ها را
دربر می‌گیرد و نسبت به ترکیب، بازگشت اولیه و جست‌وجوی نامحدودِ
اعمال‌شده بر توابع 
\emph{منظم} بسته است.
\end{defn}

آشکار است که هر تابع بازگشتی عمومی تام است. تفاوت میان
\olref{defn:general-recursive} و 
\olref[par]{defn:recursive-fn} آن است که در تعریف دوم اجازه داریم در
میانهٔ کار از توابع بازگشتی جزئی استفاده کنیم؛ تنها شرط این است که
تابع نهایی تام باشد. پس واژهٔ «عمومی»، که یادگاری تاریخی است، نامی
نابجاست: در ظاهر، 
\olref{defn:general-recursive} از
\olref[par]{defn:recursive-fn} 
\emph{کمتر} عمومی است. اما خوشبختانه این تفاوت ظاهری است؛ هرچند دو
تعریف متفاوت‌اند، مجموعهٔ توابع بازگشتی عمومی و مجموعهٔ توابع بازگشتی
یکی‌اند.

\end{document}
